\section{Covariance controls and sampling schemes}
\label{sec:sampling}

Population gradient flow depends on the covariance and fixed scales. SGD also
depends on higher update moments and on how training observations are reused.
This section keeps population centering and scales known and fixes the intercept
at its population value. The fixed-data statements condition on the entire
observed sample, including its realized response noise.

\subsection{Matched covariance, different update moments}

Let $0 < q < 1/2$, $0 \le a \le 1$, and $J \sim \operatorname{Bernoulli}(q)$.
For $U,V \in \{-1,1\}$, specify the conditional joint probabilities by
\[
 \Prob(U = u,V = v \mid J = j) = \frac{1 + a_j uv}{4},\qquad a_1
                               = a,\quad a_0
                               = -\frac{qa}{1 - q}.
\]
All cells are nonnegative. The original binary features are $W_1 = (U + 1)/2$,
$W_2 = (V + 1)/2$, and $J$. Their means and covariance are $(1/2,1/2,q)$ and
$\diag(1/4,1/4,q(1 - q))$ for every $a$. After standard-deviation scaling, put
$z = (U,V,Z)^\top$, $Z = (J - q)/\sqrt{q(1 - q)}$. Then
\[
 H = I_3,\qquad \E[UVZ^2]
   = d_a
   := \frac{a(1 - 2q)}{1 - q},\qquad \E Z^4
   = \mu_4.
\]
Thus even all marginal feature distributions agree, while a mixed fourth moment
changes.

\begin{proposition}[A covariance-matched SGD comparison]
 \label{prop:matched}
 For the distributions above, population GD, population gradient flow, and the
 fresh-SGD mean coincide under identical scales and initialization. In the
 noiseless model, initialize with original-scale coefficient error
 $(0,0,e_0)^\top$. At the same $B$ and $\eta$, after two fresh-batch updates,
 \[
  \E_a[(\beta_{R,2} - \beta_R^*)^2] - \E_0[(\beta_{R,2} - \beta_R^*)^2] = \frac{2\eta^4d_a^2}{B^2}e_0^2.
 \]
 The comparison uses $2B$ processed observations in both populations.
\end{proposition}
\begin{proof}
 Conditional balanced signs give zero means and zero common--rare covariances.
 The choice $(1 - q)a_0 + qa_1 = 0$ gives $\E UV = 0$, while conditioning on $J$
 gives the stated $d_a$. The mean claims follow from $H = I_3$. Write
 $E_{33} = e_3e_3^\top$. Then
 \[
  \mathcal M_a(E_{33}) = \begin{pmatrix}
                          1   & d_a & 0     \\
                          d_a & 1   & 0     \\
                          0   & 0   & \mu_4
                         \end{pmatrix}.
 \]
 Since $R = (1 - \eta)I_3$, the only $a$-dependent entries of
 $\mathcal T_{a,B}(E_{33})$ are its two common off-diagonal entries
 $\eta^2d_a/B$. Applying $\mathcal T_{a,B}$ again contributes $2\eta^4d_a^2/B^2$
 to its rare diagonal entry. Multiplication by the initial scaled squared error
 $q(1 - q)e_0^2$ and division by $q(1 - q)$ prove the result.
\end{proof}

Stability can differ as well. At $q = 1/4$, $\mu_4 = 7/3$ and $d_a = 2a/3$. On
symmetric matrices, the eigenvalues of $\mathcal M_a$ are
\[
 0,\quad 2,\quad 2(1 - d_a),\quad 2(1 + d_a),\quad \frac{2 + \mu_4 \pm \sqrt{(\mu_4 - 2)^2 + 8(1 + d_a^2)}}{2}.
\]
These follow by separating the common diagonal difference, the two common--rare
off-diagonal modes, and the span of the common diagonal sum, the common
off-diagonal, and the rare diagonal. The largest eigenvalue $\lambda_a$ is the
displayed plus root, so
\[
 \eta_{\mathrm{crit},a} = \frac{2}{1 + (\lambda_a - 1)/B}.
\]
For $B = 1$ the thresholds are $12/(13 + \sqrt{73})$ at $a = 0$ and
$12/(13 + \sqrt{105})$ at $a = 1$. The common step $\eta = 8/15$ is strictly
stable for the first distribution and fails all-initializations mean-square
convergence for the second. Both means converge. This comparison isolates
information absent from the population covariance. For fixed $B$ and
$q \downarrow 0$, both thresholds remain asymptotic to $2Bq$, consistent with
the leading rare-limit result in Section~\ref{sec:sgd}.

\subsection{Sampling a fixed dataset with replacement}

For fixed observations $(z_i,y_i)_{i=1}^n$, suppose
$H_D = n^{-1}\sum_i z_i z_i^\top \succ 0$, and let $\widehat\theta_D$ be the
empirical least-squares coefficient in scaled coordinates. Define
\[
 r_i = y_i - z_i^\top\widehat\theta_D,\quad g_i
     = z_i r_i,\quad \frac1n\sum_i g_i
     = 0,\quad v_k
     = \theta_k - \widehat\theta_D,\quad R_D
     = I - \eta H_D.
\]
Each update samples $B$ indices independently and uniformly with replacement,
independently of previous index draws conditional on the dataset. Write
\[
 \mathcal M_D(U) = \frac1n\sum_i z_i z_i^\top U z_i z_i^\top,\quad \mathcal V_D(U)
                 = \mathcal M_D(U) - H_DUH_D,\quad Q_D
                 = \frac1n\sum_i g_i g_i^\top,
\]
and let $\mathcal T_{D,B}(U) = R_DUR_D + \eta^2\mathcal V_D(U)/B$. The
additional linear map
\[
 J_D(m) = \frac1n\sum_i z_i z_i^\top m g_i^\top
\]
records correlation between the sampled curvature and the residual score.

\begin{proposition}[Conditional replacement moments]
 \label{prop:replacement}
 Conditional on the dataset, let $m_k = \E_Dv_k$ and $U_k = \E_D[v_kv_k^\top]$,
 with finite initial second moment and initialization independent of the sampled
 indices. Then
 \begin{align*}
  m_{k+1} & = R_Dm_k,                                                                     \\
  U_{k+1} & = \mathcal T_{D,B}(U_k) + \frac{\eta^2}{B}\{Q_D - J_D(m_k) - J_D(m_k)^\top\}.
 \end{align*}
 For $\Sigma_k = U_k - m_km_k^\top$, the covariance recurrence has the
 additional term $\eta^2\mathcal V_D(m_km_k^\top)/B$ on the right.
\end{proposition}
\begin{proof}
 Put $\widehat H = B^{-1}\sum_b z_{i_b}z_{i_b}^\top$ and
 $\bar g = B^{-1}\sum_b g_{i_b}$. The update is
 $v_{k+1} = (I - \eta\widehat H)v_k + \eta\bar g$. The normal equations give
 $\E_D\bar g = 0$ and $\E_D\bar g\bar g^\top = Q_D/B$. Independence within a
 batch and $n^{-1}\sum_i g_i = 0$ give
 $\E_D[\widehat Hm\bar g^\top] = J_D(m)/B$. Expanding the update outer product
 proves the second-moment formula. The conditional mean and covariance follow
 directly.
\end{proof}

Thus empirical normal equations do not replace the fresh-sample conditional
noise-mean assumption. The $J_D$ terms can be nonzero: scalar observations
$z_1 = 1,z_2 = 2$ with residuals $r_1 = 2,r_2 = -1$ have zero empirical score
but $J_D(m) = -3m$. The empirical mean of these features need not vanish when
centering uses known population quantities. The conditional mean follows
empirical GD toward $\widehat\theta_D$, and the original-scale error relative to
the population signal is
\[
 \beta_k - \beta^* = S^{-1}v_k + (S^{-1}\widehat\theta_D - \beta^*).
\]
The second term persists when the same noisy observations are reused. In the
noiseless full-rank model the residuals vanish, so the homogeneous operator
criterion in Proposition~\ref{prop:matrix-stability} applies with empirical
moments. Population thresholds involving $Bq$ require additional empirical
moment control before being used conditionally on a random dataset.

\subsection{Random reshuffling}

For single-observation updates, augment the error by a constant coordinate:
\[
 w = \begin{pmatrix}
      v \\
      1
     \end{pmatrix},\qquad A_i
   = \begin{pmatrix}
      I - \eta z_i z_i^\top & \eta g_i \\
      0                     & 1
     \end{pmatrix}.
\]
For a uniform random permutation $\pi$ of the $n$ indices, put
$A_\pi = A_{\pi(n)}\cdots A_{\pi(1)}$. Independently reshuffling each epoch
gives exact conditional epoch recurrences
\[
 \E_Dw_{e+1} = \E_\pi[A_\pi]\E_Dw_e,\qquad \E_D[w_{e+1}w_{e+1}^\top]
             = \E_\pi\left[ A_\pi\E_D[w_ew_e^\top]A_\pi^\top \right].
\]
Here $e$ counts epochs, each consisting of $n$ updates and $n$ processed
observations. These identities follow by conditioning on the epoch's starting
iterate and using an independent permutation. They do not factor into $n$
independent single-update operators. Fixed-size disjoint minibatches are handled
by replacing each $A_i$ with its batch-average update, specifying any incomplete
final batch.

If the original dataset consists of independent identically distributed
observations, an independent permutation preserves its product distribution.
Consequently, the first pass, unconditionally over the dataset, has exactly the
fresh-sample law for an initialization independent of the dataset and known
population centering and scales. Conditioning on the dataset or making another
pass changes this conclusion. Empirical centering, estimated scales, and an
initialization fitted to the dataset also invalidate this argument.

For a concrete conditional counterexample, take a population-standardized
Bernoulli feature with prevalence $1/3$ and a noiseless dataset containing one
observation at each of its values $\sqrt2$ and $-1/\sqrt2$. The two curvatures
are $2$ and $1/2$. At a common two-observation budget,
\[
 v_2^{\mathrm{shuffle}} = (1 - 2\eta)(1 - \eta/2)v_0,\qquad \E_Dv_2^{\mathrm{replace}}
                        = (1 - 5\eta/4)^2v_0.
\]
The difference is $-9\eta^2v_0/16$. Replacement contracts in mean square exactly
for $0 < \eta < 20/17$, while reshuffling contracts per epoch for
$0 < \eta < 5/2$. At $\eta = 2$, reshuffling reaches zero in one epoch and
replacement has second-moment factor $9/2$ per update. These statements follow
by multiplying the two scalar factors and averaging their squares, respectively.
They concern this specified dataset and do not establish a general advantage for
reshuffling.

Reusing response noise also changes unconditional error. With one observed pair,
an equiprobable sign feature $z$, independent mean-zero noise of variance
$\sigma^2$, and initialization at truth, repeated use gives
\[
 u_k = \{1 - (1 - \eta)^k\}z\varepsilon,\qquad \E u_k^2
     = \sigma^2\{1 - (1 - \eta)^k\}^2.
\]
Fresh observations instead give $\eta^2\sigma^2\sum_{j=0}^{k-1}(1 - \eta)^{2j}$.
For $0 < \eta < 2$, the limits are $\sigma^2$ and $\eta\sigma^2/(2 - \eta)$,
respectively. Conditional finite-data identities also require extra
integrability before averaging over a random dataset: noiseless reuse of a
single feature gives $u_2^2 = (1 - \eta z^2)^4u_0^2$, which can have infinite
expectation when $\E z^4 < \infty$ but $\E z^8 = \infty$.
